\section{A Pfaffian cocycle detecting the trace coefficient}
\label{sec:detector}

Retain the finitely generated group $G$, idempotent, infinite-order element
$g$, and notation $J,H,Y,v_k,\lambda,z_{2m}$ of
Section~\ref{sec:idempotents}.  We construct an $H$-invariant cocycle whose
value on $z_{2m}$ is $2^{-m}\lambda$.  The construction first records the
relative powers of $g$ between two lifts of vertices of $Y$.  An alternating
Pfaffian expression will then remove the dependence on those lifts.

\subsection{Relative exponents between lifts}

\begin{lemma}\label{lem:beta}
There is a $J$-invariant function $\beta:G\times G\to\R$ such that
\begin{equation}\label{eq:beta-properties}
 \beta(h',h)=-\beta(h,h'),\qquad
 \beta(g^r h,g^s h')=\beta(h,h')+s-r
 \quad(r,s\in\Z).
\end{equation}
\end{lemma}

\begin{proof}
The diagonal $H$-action on the set of ordered pairs $Y\times Y$ is free,
because its action on the first coordinate is free.  This also applies to
diagonal pairs $(y,y)$.  For each orbit choose a representative $(y_o,y'_o)$
and lifts $(l_o,l'_o)\in G^2$.  For $(h,h')\in G^2$, there is a unique
$\bar c\in H$ carrying the appropriate representative to the projected
pair.  Choose a lift $c\in J$ and write
\[
 h=g^r c l_o,\qquad h'=g^s c l'_o.
\]
For fixed $c$, the integers $r,s$ are unique because $g$ has infinite order.
Replacing $c$ by $g^t c$ replaces them by $r-t,s-t$.  Thus
$B(h,h')=s-r$ is well defined.  Commutation with $g$ shows that $B$ is
$J$-invariant, and independently changing the two lifts gives
\[
 B(g^a h,g^b h')=B(h,h')+b-a.
\]
The antisymmetrization
\[
 \beta(h,h')=\frac{B(h,h')-B(h',h)}2
\]
has all the required properties.
\end{proof}

The function $\beta$ itself need not descend to $Y\times Y$.  In fact,
\eqref{eq:beta-properties} gives $\beta(h,g^k h)=k$ even though these two
lifts project to the same vertex.  Nor is there a difficulty if an element
$\bar c\in H$ exchanges two distinct projected vertices.  If a lift
$c\in J$ satisfies
$c h_0=g^a h_1$ and $c h_1=g^b h_0$, its invariance merely requires
\[
 2\beta(h_0,h_1)=b-a,
\]
which is precisely what antisymmetrization gives.  Real values are useful
here: for $G=\langle t\rangle$ and $g=t^2$, the value
$\beta(1,t)=1/2$ is forced.

\subsection{The cyclic Pfaffian identity}

In degree two, a lift-independent expression is
\[
 c(y_0,y_1,y_2)
 =\beta(h_1,h_2)-\beta(h_0,h_2)+\beta(h_0,h_1),
 \qquad \bar h_i=y_i.
\]
Changing $h_i$ to $g^{r_i}h_i$ adds
$(r_2-r_1)-(r_2-r_0)+(r_1-r_0)=0$, so $c$ depends only on the
projected vertices. In higher even degrees we replace $\beta$ by a
Pfaffian and take the same alternating deletion sum. The Pfaffian
still depends on the lifts; its deletion sum will define the cocycle.
We first compute the weighted Pfaffian sum, then prove the descent
and cocycle properties of its deletion sum.

For a list $\mathbf h=(h_0,\ldots,h_{2m-1})$ with $m\ge1$, let
\[
 p_m(\mathbf h)=\operatorname{Pf}
       \bigl(\beta(h_i,h_j)\bigr)_{0\le i,j<2m},
 \qquad p_0=1.
\]
We use the convention that $p_m(\mathbf h)$ is the coefficient of
$\varepsilon_0\wedge\cdots\wedge\varepsilon_{2m-1}$ in $b^m/m!$, where
$\varepsilon_i$ are the coordinate one-forms on the space of list positions
and
\[
 b=\sum_{i<j}\beta(h_i,h_j)\,
              \varepsilon_i\wedge\varepsilon_j.
\]
In particular, $p_1(h_0,h_1)=\beta(h_0,h_1)$, and $p_m$ is alternating in
its entries and $J$-invariant.  The finite weighted sums
\[
 q_m=v_{2m}(p_m)\quad(m\ge1),\qquad q_0=\lambda
\]
are therefore defined.  No boundedness of $p_m$ on all lists is needed.

\begin{lemma}\label{lem:pfaffian-recurrence}
For every $m\ge1$, one has $q_m=\tfrac12q_{m-1}$ and hence
$q_m=2^{-m}\lambda$.
\end{lemma}

\begin{proof}
An ordinary cyclic rotation of a list of length $2m$ changes its Pfaffian
by the sign $(-1)^{2m-1}=-1$.  In the twisted rotation
\[
 T\mathbf h=(h_1,\ldots,h_{2m-1},g h_0),
\]
replacing the last entry $h_0$ by $g h_0$ adds $1$ to every matrix entry in
the last column above the diagonal.  Expansion along that column gives
\begin{equation}\label{eq:pfaffian-rotation}
 \begin{split}
 p_m(T\mathbf h)
   ={}&-p_m(\mathbf h)\\
     &+\sum_{i=1}^{2m-1}(-1)^{i+1}
       p_{m-1}(h_1,\ldots,\widehat{h_i},\ldots,h_{2m-1}).
 \end{split}
\end{equation}
Indeed the last-column cofactor in position $i$, with positions numbered
$1,\ldots,2m$, has sign $(-1)^{i+1}$; each Pfaffian matching uses exactly
one entry from that column.

For $m\ge2$, each minor on the second line omits $h_0$ and $h_i$ from the
original list.  Applying the deletion identity in
Lemma~\ref{lem:path-identities} twice shows that its weighted sum under
$v_{2m}$ is $q_{m-1}$.  For $m=1$, the minor is $p_0=1$, whose weighted sum
is the total weight $v_2(1)=\lambda=q_0$.  Thus this endpoint uses no vector
$v_0$.  The invariance of $v_{2m}$ under $T$ and
$\sum_{i=1}^{2m-1}(-1)^{i+1}=1$ now turn
\eqref{eq:pfaffian-rotation} into
\[
 q_m=-q_m+q_{m-1}.
\]
The claimed formulas follow.
\end{proof}

\subsection{Descent to an exterior cocycle}

The preceding calculation detects the coefficient on lists of lifts.
We next descend it to the exterior complex on $Y$, where the sparse-chain
theorem applies.

\begin{proposition}\label{prop:detector}
For every integer $m\ge1$, the formula
\begin{equation}\label{eq:detector-formula}
 \psi_m(y_0,\ldots,y_{2m})
   =\sum_{i=0}^{2m}(-1)^i
       p_m(h_0,\ldots,\widehat{h_i},\ldots,h_{2m}),
\end{equation}
where $h_i\in G$ is any lift of $y_i\in Y$, defines an $H$-invariant
linear functional $\psi_m:E_{2m}\to\C$ satisfying
\[
 \psi_m\partial=0,\qquad
 \psi_m(z_{2m})=2^{-m}\lambda.
\]
Here evaluation on the coinvariant chain $z_{2m}$ is defined by
$H$-invariance.
\end{proposition}

\begin{proof}
On the $2m+1$ list positions put
\[
 u=\sum_{i=0}^{2m}\varepsilon_i,
 \qquad
 b=\sum_{0\le i<j\le2m}\beta(h_i,h_j)\,
                         \varepsilon_i\wedge\varepsilon_j.
\]
The right side of \eqref{eq:detector-formula} is the coefficient of the
ordered top form in $u\wedge b^m/m!$.  Moving the factor
$\varepsilon_i$ past its $i$ preceding positions gives exactly the sign
$(-1)^i$ in that formula.

Changing the lifts to $g^{r_i}h_i$ replaces $\beta(h_i,h_j)$ by
$\beta(h_i,h_j)+r_j-r_i$.  Thus it replaces $b$ by
$b+u\wedge v$, with $v=\sum_i r_i\varepsilon_i$.  Every correction term
in $u\wedge(b+u\wedge v)^m$ contains two factors $u$ and vanishes.
This proves independence of the lifts.

The expression is alternating in the lifted list.  After lift independence
has been established, it is therefore alternating in the projected
vertices as well.  In particular, if two projected vertices coincide,
choose equal lifts for them to see that the expression vanishes.  It
consequently defines a linear functional on
$E_{2m}=\bigwedge^{2m+1}\C^{(Y)}$.
For $\bar c\in H$, choose $c\in J$ and use $c h_i$ as the lifts of the
translated vertices.  The $J$-invariance of $\beta$ proves the
$H$-invariance of $\psi_m$.

To check the cocycle identity, choose lifts simultaneously for a list of
$2m+2$ vertices.  On every face, \eqref{eq:detector-formula} is the
alternating deletion coboundary of $p_m$ on these same lifts.  A second
deletion cancels each twice-omitted list in its two possible orders.
Hence $\psi_m\partial=0$.

Finally, $z_{2m}$ is the exterior projection of $v_{2m+1}$.
Each one-deletion term in \eqref{eq:detector-formula} has weighted sum
$v_{2m}(p_m)=q_m$, by Lemma~\ref{lem:path-identities}.  Since
$\sum_{i=0}^{2m}(-1)^i=1$, their signed sum gives
\[
 \psi_m(z_{2m})=q_m=2^{-m}\lambda,
\]
using Lemma~\ref{lem:pfaffian-recurrence}.  Repeated projected vertices
already have zero exterior image and zero value under $\psi_m$, so they
require no separate restriction in this evaluation.
\end{proof}
